\newpage

\section{\huge \textcolor{blue}{一阶常系数线性方程}}

补充完复数相关知识后, 继续回到一阶常系数线性方程
{\Large  \[ y' + k y = q(t) \] }
下引入新概念:
\subsection{\large 输入-响应}
\subsubsection{\large 概念}
如上公式:\\
输入: $q(t)$\\
响应: $y(t) $ 的解\\
以前文出现过的温度传导模型为例:

\begin{figure}[htbp]
    \centering
\includegraphics[width=0.6\textwidth]{C:/Users/37/Desktop/Notes on Common Calculus Equations/Notes and pictures/cwffc7.1.1.png}
    \caption{传导模型}
    \label{fig:cwffc7.1.1}
\end{figure}

$T_e$ 外部温度为输入, $T$ 内部温度变化受外部温度影响即为响应

\subsubsection{\large 叠加原理}
\begin{center}
    \begin{tabular}{c}
        $q_1(t) \longrightarrow  y_1(t)$ \\ 
        $q_2(t) \longrightarrow  y_2(t)$   \\ \hline
        $q_1(t) + q_2(t) \longrightarrow  y_1(t) + y_2(t)$ \\
        $c q_1(t) \longrightarrow  c y_1(t)$
    \end{tabular}
\end{center}
显然输入和响应是线性的
\subsection{\large 三角实例}
下以输入为三角函数为例:
\[ y' + k y = k q_e(t) \]
对物理输入 $q_e = \cos(\omega t) $, 求响应\\
复化微积分:
\[ e^{i\omega t} = \cos(\omega t) + i \sin(\omega t) \]
\[\tilde{y'} + k \tilde{y} = k e^{i\omega t}\]
其中 $\tilde{y}$ 包含复解, $\tilde{y} = y_1 +i y_2$\\
现问题转化为:\\
求 $\tilde{y}$, 取实部$y_1$, 即为原方程解\\
积分因子: $e^{k t}$\\
解得:
\[  \tilde{y} = \frac{1}{1 + i \frac{\omega }{k}} e^{i\omega t} \]
接下来取实部(有2种方法):
\subsubsection{\large 极坐标转化}
补充:\\
若 $\alpha $ 为复数, 有下:
\begin{enumerate}
    \item  $\left\lvert \frac{1} {\alpha} \right\rvert  = \frac{1}{\left\lvert \alpha \right\rvert } $  
    \item  $\arg(\frac{1} {\alpha}) + \arg (\alpha ) = 0 $
\end{enumerate}

由上2条, 即可求出:
\[ \frac{1}{1 + i \frac{\omega }{k}} = \frac{1}{\sqrt{1 + (\frac{\omega }{k})^2} } e^{- i \phi } \]
其中 $\arg (1 + i \frac{\omega }{k}) = \phi $
\[ \therefore \hspace*{2em} \tilde{y} = \frac{1}{\sqrt{1 + (\frac{\omega}{k})^2}} e^{i(\omega t - \phi)} \]
\[ y_1 = Re(\tilde{y}) = \frac{1}{\sqrt{1 + (\frac{\omega}{k})^2}} \cos(\omega t - \phi) \]
其中$\phi = \arctan (\frac{\omega }{k })$


\subsubsection{\large 笛卡尔坐标转化}
\[ e^{i\omega t} = \cos(\omega t) + i \sin(\omega t) \]
\[ \tilde{y} = \frac{1 - i \frac{\omega }{k}}{1 + (\frac{\omega }{k})^2} (\cos(\omega t) + i \sin(\omega t)) \]
\[ y_1 = Re(\tilde{y}) = \frac{1}{1 + (\frac{\omega }{k})^2} (\cos(\omega t) + \frac{\omega }{k} \sin(\omega t))\]

通过2种转化法得2种答案形式, 在正常观念中, 2值应相等
, 事实也正是如此\\
联系:\\
显然, 发现符合三角函数里的辅助角公式:
\[ a \cos\theta  + b \sin\theta = A \cos(\theta - \phi ) \]
其中 $A = \sqrt{a^2 + b^2}$, $\phi = \arctan\frac{b}{a}$
\subsection{\large 基本线性方程}
\begin{enumerate}
    \item $y' + k y = k q_e(t)$, $k > 0$
    \item $y' + k y = q(t)$, $k > 0$
    \item $y' + p(t) y = q(t)$
\end{enumerate}
以上是前文所涉及到的基本线性方程, 其$k$ 全保持大于$0$, 
若$k < 0$, 前文所涉及的"暂态、稳态、输入、响应"概念不再适用, 
但方程的解法不改变, 简单说即为解发散.